\chapter{Einzelschritte des Beweises}

Die Argumente \textup{(M1)} bis \textup{(M17)} führen die einzelnen
Schritte der Konstruktion vollständig aus. Die Definitionen, Aussagen
und Beweise dieses Kapitels gehören zu dieser Ergänzung; der Hauptband
enthält den daraus folgenden Existenzsatz.
Den zusammenhängenden Beweis enthält die \MogReadingLink{}.

\section{Träger und Grundhalbgruppen}



\MogStatementTarget{MogiljanskajaRuleDef}
\FormulaDefDeltaK[Die Fallrelation der beiden Multiplikationen]{%
  \begin{aligned}[t]
    \MogRule{x}{y}{z}\ \leftrightarrow{}&
      \exists i,j\in\mathbb N\,
        (x=a_i\land y=a_j\land z=b_{ij})\ \lor{}\\[-2pt]
      &((x\notin L\lor y\notin L)\land z=\mathbf0)
  \end{aligned}
}{MogiljanskajaRuleDef}{%
  \DeltaRow{Mengen}{L\dsep\mathbf0\dsep x\dsep y\dsep z
    \dsep a_i\dsep a_j\dsep b_{ij}\dsep i\dsep j}
  \DeltaRow{Prädikatensymbole}{\mathsf R_{\mathrm M}}
}

\MogStatementTarget{MogiljanskajaArgumentM1WellDefined}
\FormulaThmDeltaKR[Argument (M1): Wohldefiniertheit der Konstruktion]{%
 \begin{aligned}[t]
  &\text{(i)}\quad L\cap D=L\cap\{\mathbf0\}=D\cap\{\mathbf0\}\\[-2pt]
  &\qquad =L\cap\{k\}=D\cap\{k\}=\{\mathbf0\}\cap\{k\}=\varnothing\\[-2pt]
  &\text{(ii)}\quad s\in L\vdash\exists!i\in\mathbb N\,(s=a_i)\\[-2pt]
  &\text{(iii)}\quad x,y\in S\vdash\exists!z\in S\,\MogRule{x}{y}{z}\\[-2pt]
  &\text{(iv)}\quad x,y\in T\vdash\exists!z\in T\,\MogRule{x}{y}{z}\\[-2pt]
  &\text{(v)}\quad \MogStar\colon S\times S\to S\\[-2pt]
  &\text{(vi)}\quad \MogDiamond\colon T\times T\to T\\[-2pt]
  &\text{(vii)}\quad x,y\in S\vdash\MogRule{x}{y}{x\MogStar y}\\[-2pt]
  &\text{(viii)}\quad x,y\in T\vdash\MogRule{x}{y}{x\MogDiamond y}
 \end{aligned}%
}{%
  \begin{aligned}[t]
    &L\cap D=L\cap\{\mathbf0\}=D\cap\{\mathbf0\}
      =L\cap\{k\}=D\cap\{k\}=\{\mathbf0\}\cap\{k\}
      =\varnothing,\\[-2pt]
    &s\in L\vdash\exists!i\in\mathbb N\,(s=a_i),\\[-2pt]
    &x,y\in S\vdash\exists!z\in S\,\MogRule{x}{y}{z},\\[-2pt]
    &x,y\in T\vdash\exists!z\in T\,\MogRule{x}{y}{z},\\[-2pt]
    &\MogStar\colon S\times S\to S\dsep
      \MogDiamond\colon T\times T\to T,\\[-2pt]
    &x,y\in S\vdash\MogRule{x}{y}{x\MogStar y},\\[-2pt]
    &x,y\in T\vdash\MogRule{x}{y}{x\MogDiamond y}
  \end{aligned}
}{MogiljanskajaArgumentM1WellDefined}{%
  \DeltaRow{Mengen}{L\dsep D\dsep S\dsep T\dsep\mathbf0\dsep k
    \dsep s\dsep x\dsep y\dsep z\dsep a_i\dsep a_j\dsep b_{ij}
    \dsep i\dsep j}
  \DeltaRow{Funktionen}{a\dsep\MogStar\dsep\MogDiamond}
}
\MogProof{MogiljanskajaArgumentM1WellDefined}{%
\begin{tabproofwide}
  \proofstepwidestar[]{L\cap D=\varnothing}{%
    \FormulaRefAuto{MogiljanskajaLayersDisjoint}[thm]}
  \proofstepwidestar[]{\mathbf0\notin L\cup D}{%
    \text{\parbox[t]{\linewidth}{\raggedright \normalfont Tag 2 ist weder 0 noch 1}}}
  \proofstepwidestarbreak[]{k\notin L\cup D\cup\{\mathbf0\}}{%
    \text{\parbox[t]{\linewidth}{\raggedright \normalfont Tag 3 ist von 0, 1 und 2 verschieden}}}
  \proofstepwidestarbreak[1,2,3]{%
    \begin{aligned}[t]
      L\cap D&=L\cap\{\mathbf0\}=D\cap\{\mathbf0\}\\[-2pt]
        &=L\cap\{k\}=D\cap\{k\}
          =\{\mathbf0\}\cap\{k\}=\varnothing
    \end{aligned}}{%
    \text{\parbox[t]{\linewidth}{\raggedright \normalfont paarweise Disjunktheit}}}
  \proofstepwidestarbreak[]{s\in L\leftrightarrow
    \exists i\in\mathbb N\,(s=a_i)}{%
    \FormulaRefAuto{MogiljanskajaLayerLMembership}[thm]}
  \proofstepwidestarbreak[5]{s\in L\vdash
    \exists!i\in\mathbb N\,(s=a_i)}{%
    \FormulaRefAuto{MogiljanskajaAIndexUnique}[thm]}
  \proofstepwidestarbreak[]{i,j\in\mathbb N\vdash
    b_{ij}\in D\subseteq S\subseteq T}{%
    \FormulaRefAuto{MogiljanskajaBElementsDef}[def],
    \FormulaRefAuto{MogiljanskajaPairDef}[def]}
  \proofstepwidestar[]{\mathbf0\in S\subseteq T}{%
    \FormulaRefAuto{MogiljanskajaPairDef}[def]}
  \proofstepwidestar[]{i,j,m,n\in\mathbb N\dsep b_{ij}=b_{mn}
    \vdash i=m\land j=n}{%
    \FormulaRefAuto{MogiljanskajaBIndexUnique}[thm]}
  \proofstepwidestarbreak[4,6,7,8,9]{x,y\in S\vdash
    \exists!z\in S\,\MogRule{x}{y}{z}}{%
    \FormulaRefAuto{MogiljanskajaRuleDef}[def]}
  \proofstepwidestarbreak[4,6,7,8,9]{x,y\in T\vdash
    \exists!z\in T\,\MogRule{x}{y}{z}}{%
    \FormulaRefAuto{MogiljanskajaRuleDef}[def]}
  \proofstepwidestarbreak[]{\MogStar\colon S\times S\to S\dsep
    \MogDiamond\colon T\times T\to T}{%
    \FormulaRefAuto{MogiljanskajaPairDef}[def]}
  \proofstepwidestarbreak[10,12]{x,y\in S\vdash
    \MogRule{x}{y}{x\MogStar y}}{%
    \FormulaRefAuto{MogiljanskajaPairDef}[def]}
  \proofstepwidestarbreak[11,12]{x,y\in T\vdash
    \MogRule{x}{y}{x\MogDiamond y}}{%
    \FormulaRefAuto{MogiljanskajaPairDef}[def]}
\end{tabproofwide}
}


\MogStatementTarget{MogiljanskajaArgumentM2Semigroups}
\FormulaThmDeltaKR[Argument (M2): Produktbereich und Nullabsorption]{%
 \begin{aligned}[t]
  &\text{(i)}\quad x,y\in S\vdash x\MogStar y\in D\cup\{\mathbf0\}\\[-2pt]
  &\text{(ii)}\quad u\in D\cup\{\mathbf0\}\dsep z\in S \vdash u\MogStar z=\mathbf0=z\MogStar u\\[-2pt]
  &\text{(iii)}\quad x,y,z\in S\vdash(x\MogStar y)\MogStar z=\mathbf0 =x\MogStar(y\MogStar z)\\[-2pt]
  &\text{(iv)}\quad x,y\in T\vdash x\MogDiamond y\in D\cup\{\mathbf0\}\\[-2pt]
  &\text{(v)}\quad u\in D\cup\{\mathbf0\}\dsep z\in T \vdash u\MogDiamond z=\mathbf0=z\MogDiamond u\\[-2pt]
  &\text{(vi)}\quad x,y,z\in T\vdash(x\MogDiamond y)\MogDiamond z=\mathbf0 =x\MogDiamond(y\MogDiamond z)\\[-2pt]
  &\text{(vii)}\quad \SemigroupAlg{S}{\MogStar}\land \SemigroupAlg{T}{\MogDiamond}
 \end{aligned}%
}{%
  \begin{aligned}[t]
    x,y\in S&\vdash x\MogStar y\in D\cup\{\mathbf0\},\\[-2pt]
    u\in D\cup\{\mathbf0\}\dsep z\in S
      &\vdash u\MogStar z=\mathbf0=z\MogStar u,\\[-2pt]
    x,y,z\in S&\vdash
      (x\MogStar y)\MogStar z=\mathbf0
        =x\MogStar(y\MogStar z),\\[2pt]
    x,y\in T&\vdash x\MogDiamond y\in D\cup\{\mathbf0\},\\[-2pt]
    u\in D\cup\{\mathbf0\}\dsep z\in T
      &\vdash u\MogDiamond z=\mathbf0=z\MogDiamond u,\\[-2pt]
    x,y,z\in T&\vdash
      (x\MogDiamond y)\MogDiamond z=\mathbf0
        =x\MogDiamond(y\MogDiamond z),\\[2pt]
    &\SemigroupAlg{S}{\MogStar}\land
      \SemigroupAlg{T}{\MogDiamond}
  \end{aligned}
}{MogiljanskajaArgumentM2Semigroups}{%
  \DeltaRow{Mengen}{L\dsep D\dsep S\dsep T\dsep\mathbf0
    \dsep x\dsep y\dsep z\dsep u}
  \DeltaRow{Funktionen}{\MogStar\dsep\MogDiamond}
}
\MogProof{MogiljanskajaArgumentM2Semigroups}{%
\begin{tabproofwide}
  \proofstepwidestar[]{x,y\in S\vdash
    x\MogStar y\in D\cup\{\mathbf0\}}{%
    \FormulaRefAuto{MogiljanskajaPairDef}[def]}
  \proofstepwidestar[]{x,y\in T\vdash
    x\MogDiamond y\in D\cup\{\mathbf0\}}{%
    \FormulaRefAuto{MogiljanskajaPairDef}[def]}
  \proofstepwidestar[]{(D\cup\{\mathbf0\})\cap L=\varnothing}{%
    \FormulaRefAuto{MogiljanskajaArgumentM1WellDefined}[thm]}
  \proofstepwidestarbreak[3]{u\in D\cup\{\mathbf0\}\dsep z\in S
    \vdash u\MogStar z=\mathbf0=z\MogStar u}{%
    \FormulaRefAuto{MogiljanskajaPairDef}[def]}
  \proofstepwidestarbreak[3]{u\in D\cup\{\mathbf0\}\dsep z\in T
    \vdash u\MogDiamond z=\mathbf0=z\MogDiamond u}{%
    \FormulaRefAuto{MogiljanskajaPairDef}[def]}
  \proofstepwidestarbreak[1,4]{x,y,z\in S\vdash
    (x\MogStar y)\MogStar z=\mathbf0
      =x\MogStar(y\MogStar z)}{%
    \text{\parbox[t]{\linewidth}{\raggedright \normalfont zweimal Nullabsorption}}}
  \proofstepwidestarbreak[2,5]{x,y,z\in T\vdash
    (x\MogDiamond y)\MogDiamond z=\mathbf0
      =x\MogDiamond(y\MogDiamond z)}{%
    \text{\parbox[t]{\linewidth}{\raggedright \normalfont zweimal Nullabsorption}}}
  \proofstepwidestar[6,7]{%
    \SemigroupAlg{S}{\MogStar}\land
    \SemigroupAlg{T}{\MogDiamond}}{%
    \FormulaRefAuto{SemigroupDef}[def]}
\end{tabproofwide}
}


\MogStatementTarget{MogiljanskajaArgumentM3Infinite}
\FormulaThmDeltaK[Argument (M3): Unendlichkeit beider Träger]{%
  \neg\Finite{S}\land\neg\Finite{T}%
}{MogiljanskajaArgumentM3Infinite}{%
  \DeltaRow{Mengen}{L\dsep S\dsep T}
  \DeltaRow{Funktionen}{a\dsep H}
}
\MogProof{MogiljanskajaArgumentM3Infinite}{%
\begin{tabproofwide}
  \proofstepwidestar[]{a\colon\mathbb N\inj L}{%
    \FormulaRefAuto{MogiljanskajaAElementsDef}[def],
    \FormulaRefAuto{MogiljanskajaAIndexUnique}[thm]}
  \proofstepwidestar[]{L\subseteq S\subseteq T}{%
    \FormulaRefAuto{MogiljanskajaPairDef}[def]}
  \proofstepwidestar[1,2]{a\colon\mathbb N\inj S}{%
    \FormulaRefAuto{InjectiveCodomainExtension}[thm]}
  \proofstepwidestar[1,2]{a\colon\mathbb N\inj T}{%
    \FormulaRefAuto{InjectiveCodomainExtension}[thm]}
  \proofstepwidestar[]{\Finite{S}\vdash
    \neg\exists H\colon\mathbb N\inj S}{%
    \FormulaRefAuto{FiniteNoNatInjection}[thm]}
  \proofstepwidestar[3,5]{\neg\Finite{S}}{\text{\parbox[t]{\linewidth}{\raggedright %
    Widerspruch mit dem Zeugen \(a\)}}}
  \proofstepwidestar[]{\Finite{T}\vdash
    \neg\exists H\colon\mathbb N\inj T}{%
    \FormulaRefAuto{FiniteNoNatInjection}[thm]}
  \proofstepwidestar[4,7]{\neg\Finite{T}}{\text{\parbox[t]{\linewidth}{\raggedright %
    Widerspruch mit dem Zeugen \(a\)}}}
  \proofstepwidestar[6,8]{\neg\Finite{S}\land\neg\Finite{T}}{\rAI{6,8}}
\end{tabproofwide}
}


\MogStatementTarget{MogiljanskajaArgumentM4ProductStatus}
\FormulaThmDeltaKR[Argument (M4): Produktstatus der ausgezeichneten Elemente]{%
 \begin{aligned}[t]
  &\text{(i)}\quad s\in S\setminus L \vdash\exists x\in S\,\exists y\in S\,(s=x\MogStar y)\\[-2pt]
  &\text{(ii)}\quad x,y\in T\vdash x\MogDiamond y\neq k
 \end{aligned}%
}{%
  \begin{aligned}[t]
    s\in S\setminus L
      &\vdash\exists x\in S\,\exists y\in S\,(s=x\MogStar y),\\[-2pt]
    x,y\in T&\vdash x\MogDiamond y\neq k
  \end{aligned}
}{MogiljanskajaArgumentM4ProductStatus}{%
  \DeltaRow{Mengen}{L\dsep D\dsep S\dsep T\dsep\mathbf0\dsep k
    \dsep s\dsep x\dsep y\dsep a_i\dsep a_j\dsep b_{ij}
    \dsep i\dsep j}
  \DeltaRow{Funktionen}{\MogStar\dsep\MogDiamond}
}
\MogProof{MogiljanskajaArgumentM4ProductStatus}{%
\begin{tabproofwide}
  \proofstepwidestar[]{s\in S\setminus L\vdash
    s\in D\lor s=\mathbf0}{%
    \FormulaRefAuto{MogiljanskajaPairDef}[def],
    \FormulaRefAuto{MogiljanskajaArgumentM1WellDefined}[thm]}
  \proofstepwidestar[]{s\in D\vdash
    \exists i,j\in\mathbb N\,(s=b_{ij})}{%
    \FormulaRefAuto{MogiljanskajaLayerDMembership}[thm]}
  \proofstepwidestar[2]{s\in D\vdash
    \exists x\in S\,\exists y\in S\,(s=x\MogStar y)}{\text{\parbox[t]{\linewidth}{\raggedright %
    Wähle \(x=a_i\) und \(y=a_j\)}}}
  \proofstepwidestar[]{\mathbf0=\mathbf0\MogStar\mathbf0}{%
    \FormulaRefAuto{MogiljanskajaPairDef}[def]}
  \proofstepwidestarbreak[1,3,4]{s\in S\setminus L\vdash
    \exists x\in S\,\exists y\in S\,(s=x\MogStar y)}{%
    \text{\parbox[t]{\linewidth}{\raggedright \normalfont Fallunterscheidung}}}
  \proofstepwidestar[]{x,y\in T\vdash
    x\MogDiamond y\in D\cup\{\mathbf0\}}{%
    \FormulaRefAuto{MogiljanskajaArgumentM2Semigroups}[thm]}
  \proofstepwidestar[]{k\notin D\cup\{\mathbf0\}}{%
    \FormulaRefAuto{MogiljanskajaArgumentM1WellDefined}[thm]}
  \proofstepwidestar[6,7]{x,y\in T\vdash x\MogDiamond y\neq k}{%
    \text{\parbox[t]{\linewidth}{\raggedright \normalfont Ausschluss aus dem Produktbereich}}}
\end{tabproofwide}
}


\MogStatementTarget{MogiljanskajaArgumentM5NotIsomorphic}
\FormulaThmDeltaK[Argument (M5): Kollaps eines angenommenen Isomorphismus]{%
  \neg\exists f\,\SemigroupIso{f}{S,\MogStar}{T,\MogDiamond}%
}{MogiljanskajaArgumentM5NotIsomorphic}{%
  \DeltaRow{Mengen}{L\dsep D\dsep S\dsep T\dsep\mathbf0\dsep k}
  \DeltaRow{Mengen}{s\dsep p\dsep q\dsep u\dsep v\dsep x\dsep y}
  \DeltaRow{Mengen}{a_i\dsep a_m\dsep a_n\dsep a_0\dsep a_1}
  \DeltaRow{Mengen}{b_{mn}\dsep b_{i0}\dsep b_{i1}\dsep i\dsep m\dsep n}
  \DeltaRow{Funktionen}{f\dsep\MogStar\dsep\MogDiamond}
}
\MogProof{MogiljanskajaArgumentM5NotIsomorphic}{%
\begin{tabproofwide}
  \proofstepwidestar[1]{\exists f\,
    \SemigroupIso{f}{S,\MogStar}{T,\MogDiamond}}{\rA}
  \proofstepwidestar[2]{\SemigroupIso{f}{S,\MogStar}{T,\MogDiamond}}{\rA}
  \proofstepwidestar[2]{f\colon S\bij T}{%
    \FormulaRefAuto{SemigroupIsoDef}[def]{2}}
  \proofstepwidestar[2]{x,y\in S\vdash
    f(x\MogStar y)=f(x)\MogDiamond f(y)}{%
    \FormulaRefAuto{SemigroupIsoDef}[def]{2}}
  \proofstepwidestar[3]{\exists s\in S\,(f(s)=k)}{\text{\parbox[t]{\linewidth}{\raggedright %
    Surjektivität und \(k\in T\)}}}
  \proofstepwidestar[5]{s\in S\dsep f(s)=k}{\rA}
  \proofstepwidestar[6]{s\notin L\vdash
    s\in D\lor s=\mathbf0}{%
    \FormulaRefAuto{MogiljanskajaPairDef}[def],
    \FormulaRefAuto{MogiljanskajaArgumentM1WellDefined}[thm]}
  \MogProofStepWideStarKeep[7]{s\notin L\vdash
    \exists p\in S\,\exists q\in S\,(s=p\MogStar q)}{\text{\parbox[t]{\linewidth}{\raggedright \FormulaRefAuto{MogiljanskajaArgumentM4ProductStatus}[thm] ; \normalfont Zeugen im \ensuremath{D} \normalfont{-Fall: } \ensuremath{a_m,\allowbreak a_n
    \dsep} \normalfont im Nullfall: \ensuremath{\mathbf0,\allowbreak \mathbf0}}}}
  \MogProofStepWideStarKeep[3,4,6,8]{s\notin L\vdash
    \exists u\in T\,\exists v\in T\,(k=u\MogDiamond v)}{\text{\parbox[t]{\linewidth}{\raggedright \normalfont Wähle \ensuremath{u=\allowbreak f(p)} \normalfont{ und } \ensuremath{v=\allowbreak f(q)}}}}
  \proofstepwidestar[]{%
    \neg\exists u\in T\,\exists v\in T\,(k=u\MogDiamond v)}{%
    \FormulaRefAuto{MogiljanskajaArgumentM4ProductStatus}[thm]}
  \proofstepwidestar[9,10]{s\in L}{\text{\parbox[t]{\linewidth}{\raggedright %
    Ausschluss der Alternative \(s\notin L\)}}}
  \proofstepwidestar[11]{\exists i\in\mathbb N\,(s=a_i)}{%
    \FormulaRefAuto{MogiljanskajaLayerLMembership}[thm]{11}}
  \proofstepwidestar[12]{i\in\mathbb N\dsep s=a_i}{\rA}
  \proofstepwidestar[4,6,13]{f(b_{i0})=\mathbf0}{%
    \ensuremath{b_{i0}=a_i\MogStar a_0\dsep
      k\MogDiamond f(a_0)=\mathbf0}}
  \proofstepwidestar[4,6,13]{f(b_{i1})=\mathbf0}{%
    \ensuremath{b_{i1}=a_i\MogStar a_1\dsep
      k\MogDiamond f(a_1)=\mathbf0}}
  \proofstepwidestar[13]{b_{i0}\neq b_{i1}}{%
    \FormulaRefAuto{MogiljanskajaBIndexUnique}[thm],
    \FormulaRefAuto{PeanoZeroNeOne}[thm]}
  \proofstepwidestar[3,14,15,16]{\bot}{\text{\parbox[t]{\linewidth}{\raggedright %
    Injektivität von \(f\)}}}
  \proofstepwidestar[2,6]{\bot}{\rEE{12,13,17}}
  \proofstepwidestar[2]{\bot}{\rEE{5,6,18}}
  \proofstepwidestar[1]{\bot}{\rEE{1,2,19}}
  \proofstepwidestar[]{\neg\exists f\,
    \SemigroupIso{f}{S,\MogStar}{T,\MogDiamond}}{\rCI{1,20}}
\end{tabproofwide}
}


\section{Mengenprodukte und Reserve}

\MogStatementTarget{MogiljanskajaArgumentM6ProductClassification}
\FormulaThmDeltaKR[Argument (M6): Elementweise Produktklassifikation]{%
 \begin{aligned}[t]
  &\text{(i)}\quad X,Y\in\PowNE{S}\vdash{}\\[-2pt]
  &\qquad\bigl(z\in X\MogStar_{\mathcal P}Y \leftrightarrow z\in R(X,Y)\cup\varepsilon(X,Y)\bigr)\\[-2pt]
  &\text{(ii)}\quad X,Y\in\PowNE{T}\vdash{}\\[-2pt]
  &\qquad\bigl(z\in X\MogDiamond_{\mathcal P}Y \leftrightarrow z\in R(X,Y)\cup\varepsilon(X,Y)\bigr)\\[-2pt]
  &\text{(iii)}\quad X,Y\in\PowNE{S}\vdash X\MogStar_{\mathcal P}Y=R(X,Y)\cup\varepsilon(X,Y)\\[-2pt]
  &\text{(iv)}\quad X,Y\in\PowNE{T}\vdash X\MogDiamond_{\mathcal P}Y=R(X,Y)\cup\varepsilon(X,Y)
 \end{aligned}%
}{%
  \begin{aligned}[t]
    X,Y\in\PowNE{S}&\vdash
      \bigl(z\in X\MogStar_{\mathcal P}Y
        \leftrightarrow z\in R(X,Y)\cup\varepsilon(X,Y)\bigr),\\[-2pt]
    X,Y\in\PowNE{T}&\vdash
      \bigl(z\in X\MogDiamond_{\mathcal P}Y
        \leftrightarrow z\in R(X,Y)\cup\varepsilon(X,Y)\bigr),\\[2pt]
    X,Y\in\PowNE{S}&\vdash
      X\MogStar_{\mathcal P}Y=R(X,Y)\cup\varepsilon(X,Y),\\[-2pt]
    X,Y\in\PowNE{T}&\vdash
      X\MogDiamond_{\mathcal P}Y=R(X,Y)\cup\varepsilon(X,Y)
  \end{aligned}
}{MogiljanskajaArgumentM6ProductClassification}{%
  \DeltaRow{Mengen}{L\dsep S\dsep T\dsep X\dsep Y
    \dsep R(X,Y)\dsep\varepsilon(X,Y)}
  \DeltaRow{Funktionen}{\MogStar_{\mathcal P}
    \dsep\MogDiamond_{\mathcal P}}
}
\MogProof{MogiljanskajaArgumentM6ProductClassification}{%
\begin{tabproofwide}
  \proofstepwidestarbreak[]{X,Y\in\PowNE{S}\dsep
    z\in X\MogStar_{\mathcal P}Y\vdash
    \exists x\in X\,\exists y\in Y\,(z=x\MogStar y)}{%
    \FormulaRefAuto{PowerSemigroupProductWitnesses}[thm]}
  \proofstepwidestarbreak[]{x\in X\cap L\dsep y\in Y\cap L\vdash
    x\MogStar y\in R(X,Y)}{%
    \FormulaRefAuto{MogiljanskajaLayerLMembership}[thm],
    \FormulaRefAuto{MogiljanskajaProductPartsDef}[def]}
  \proofstepwidestarbreak[]{x\in X\dsep y\in Y\dsep
    (x\notin L\lor y\notin L)\vdash
    x\MogStar y=\mathbf0\in\varepsilon(X,Y)}{%
    \FormulaRefAuto{MogiljanskajaPairDef}[def],
    \FormulaRefAuto{MogiljanskajaProductPartsDef}[def]}
  \proofstepwidestarbreak[1,2,3]{X,Y\in\PowNE{S}\dsep
    z\in X\MogStar_{\mathcal P}Y\vdash
    z\in R(X,Y)\cup\varepsilon(X,Y)}{%
    \text{\parbox[t]{\linewidth}{\raggedright \normalfont Fallunterscheidung für die Produktzeugen}}}
  \proofstepwidestarbreak[]{z\in R(X,Y)\vdash
    z\in X\MogStar_{\mathcal P}Y}{%
    \FormulaRefAuto{MogiljanskajaProductPartsDef}[def],
    \FormulaRefAuto{MogiljanskajaPairDef}[def]}
  \proofstepwidestarbreak[]{%
    \begin{aligned}[t]
      X\neq\varnothing\dsep Y\neq\varnothing\dsep
        &\neg(X\subseteq L\land Y\subseteq L)\vdash{}\\[-2pt]
      &\exists x\in X\,\exists y\in Y\,(x\notin L\lor y\notin L)
    \end{aligned}}{\text{\parbox[t]{\linewidth}{\raggedright \normalfont Fälle \ensuremath{X\nsubseteq L} \normalfont{ und } \ensuremath{Y\nsubseteq L} \normalfont{; der andere Faktor ist nichtleer}}}}
  \proofstepwidestarbreak[6]{X,Y\in\PowNE{S}\dsep
    z\in\varepsilon(X,Y)\vdash
    z\in X\MogStar_{\mathcal P}Y}{%
    \FormulaRefAuto{MogiljanskajaProductPartsDef}[def],
    \FormulaRefAuto{MogiljanskajaPairDef}[def]}
  \proofstepwidestarbreak[5,7]{X,Y\in\PowNE{S}\dsep
    z\in R(X,Y)\cup\varepsilon(X,Y)\vdash
    z\in X\MogStar_{\mathcal P}Y}{%
    \text{\parbox[t]{\linewidth}{\raggedright \normalfont Vereinigungsfall}}}
  \proofstepwidestarbreak[4,8]{X,Y\in\PowNE{S}\vdash
    \bigl(z\in X\MogStar_{\mathcal P}Y
      \leftrightarrow z\in R(X,Y)\cup\varepsilon(X,Y)\bigr)}{%
    \text{\parbox[t]{\linewidth}{\raggedright \normalfont Äquivalenzeinführung}}}
  \proofstepwidestarbreak[]{X,Y\in\PowNE{T}\dsep
    z\in X\MogDiamond_{\mathcal P}Y\vdash
    \exists x\in X\,\exists y\in Y\,(z=x\MogDiamond y)}{%
    \FormulaRefAuto{PowerSemigroupProductWitnesses}[thm]}
  \proofstepwidestarbreak[]{x\in X\cap L\dsep y\in Y\cap L\vdash
    x\MogDiamond y\in R(X,Y)}{%
    \FormulaRefAuto{MogiljanskajaLayerLMembership}[thm],
    \FormulaRefAuto{MogiljanskajaProductPartsDef}[def]}
  \proofstepwidestarbreak[]{x\in X\dsep y\in Y\dsep
    (x\notin L\lor y\notin L)\vdash
    x\MogDiamond y=\mathbf0\in\varepsilon(X,Y)}{%
    \FormulaRefAuto{MogiljanskajaPairDef}[def],
    \FormulaRefAuto{MogiljanskajaProductPartsDef}[def]}
  \proofstepwidestarbreak[10,11,12]{X,Y\in\PowNE{T}\dsep
    z\in X\MogDiamond_{\mathcal P}Y\vdash
    z\in R(X,Y)\cup\varepsilon(X,Y)}{%
    \text{\parbox[t]{\linewidth}{\raggedright \normalfont Fallunterscheidung für die Produktzeugen}}}
  \proofstepwidestarbreak[]{z\in R(X,Y)\vdash
    z\in X\MogDiamond_{\mathcal P}Y}{%
    \FormulaRefAuto{MogiljanskajaProductPartsDef}[def],
    \FormulaRefAuto{MogiljanskajaPairDef}[def]}
  \proofstepwidestarbreak[6]{X,Y\in\PowNE{T}\dsep
    z\in\varepsilon(X,Y)\vdash
    z\in X\MogDiamond_{\mathcal P}Y}{%
    \FormulaRefAuto{MogiljanskajaProductPartsDef}[def],
    \FormulaRefAuto{MogiljanskajaPairDef}[def]}
  \proofstepwidestarbreak[14,15]{X,Y\in\PowNE{T}\dsep
    z\in R(X,Y)\cup\varepsilon(X,Y)\vdash
    z\in X\MogDiamond_{\mathcal P}Y}{%
    \text{\parbox[t]{\linewidth}{\raggedright \normalfont Vereinigungsfall}}}
  \proofstepwidestarbreak[13,16]{X,Y\in\PowNE{T}\vdash
    \bigl(z\in X\MogDiamond_{\mathcal P}Y
      \leftrightarrow z\in R(X,Y)\cup\varepsilon(X,Y)\bigr)}{%
    \text{\parbox[t]{\linewidth}{\raggedright \normalfont Äquivalenzeinführung}}}
  \proofstepwidestarbreak[9]{\begin{aligned}[t]&X,Y\in\PowNE{S}\vdash{}\\[-2pt]
    &X\MogStar_{\mathcal P}Y=R(X,Y)\cup\varepsilon(X,Y)\end{aligned}}{%
    \text{\parbox[t]{\linewidth}{\raggedright \normalfont Extensionalität}}}
  \proofstepwidestarbreak[17]{\begin{aligned}[t]&X,Y\in\PowNE{T}\vdash{}\\[-2pt]
    &X\MogDiamond_{\mathcal P}Y=R(X,Y)\cup\varepsilon(X,Y)\end{aligned}}{%
    \text{\parbox[t]{\linewidth}{\raggedright \normalfont Extensionalität}}}
\end{tabproofwide}
}


\MogStatementTarget{MogiljanskajaArgumentM7AbsorptionHypotheses}
\FormulaThmDeltaKR[Argument (M7): Voraussetzungen der Reserveabsorption]{%
 \begin{aligned}[t]
  &\text{(i)}\quad P\cap U=\varnothing\\[-2pt]
  &\text{(ii)}\quad c_2\notin P\cup U\\[-2pt]
  &\text{(iii)}\quad \{k\}\cap D=\varnothing\\[-2pt]
  &\text{(iv)}\quad \sigma\colon U\bij V\\[-2pt]
  &\text{(v)}\quad \theta\colon U\bij D
 \end{aligned}%
}{%
  \begin{aligned}[t]
    &P\cap U=\varnothing\dsep c_2\notin P\cup U\dsep
      \{k\}\cap D=\varnothing,\\[-2pt]
    &\sigma\colon U\bij V\dsep\theta\colon U\bij D
  \end{aligned}
}{MogiljanskajaArgumentM7AbsorptionHypotheses}{%
  \DeltaRow{Mengen}{D\dsep U\dsep V\dsep P\dsep c_2\dsep k}
  \DeltaRow{Funktionen}{\sigma\dsep\theta}
}
\MogProof{MogiljanskajaArgumentM7AbsorptionHypotheses}{%
\begin{tabproofwide}
  \proofstepwidestar[]{P\cap U=\varnothing\dsep c_2\notin P}{%
    \FormulaRefAuto{MogiljanskajaReserveCoreFacts}[thm]}
  \proofstepwidestar[]{c_2\notin U}{%
    \FormulaRefAuto{MogiljanskajaReserveMarkerFacts}[thm]}
  \proofstepwidestar[1,2]{c_2\notin P\cup U}{%
    \text{\parbox[t]{\linewidth}{\raggedright \normalfont Vereinigungskriterium}}}
  \proofstepwidestar[]{k\notin D}{%
    \FormulaRefAuto{MogiljanskajaArgumentM1WellDefined}[thm]}
  \proofstepwidestar[4]{\{k\}\cap D=\varnothing}{%
    \text{\parbox[t]{\linewidth}{\raggedright \normalfont Einermengenkriterium}}}
  \proofstepwidestar[]{\sigma\colon U\bij V}{%
    \FormulaRefAuto{MogiljanskajaSigmaBijective}[thm]}
  \proofstepwidestar[]{\theta\colon U\bij D}{%
    \FormulaRefAuto{MogiljanskajaThetaBijective}[thm]}
  \proofstepwidestarbreak[1,3,5,6,7]{%
    \begin{aligned}[t]
      &P\cap U=\varnothing\dsep c_2\notin P\cup U\dsep
        \{k\}\cap D=\varnothing,\\[-2pt]
      &\sigma\colon U\bij V\dsep\theta\colon U\bij D
    \end{aligned}}{\text{\parbox[t]{\linewidth}{\raggedright \normalfont Bündelung}}}
\end{tabproofwide}
}


\pagebreak[2]
\MogStatementTarget{MogiljanskajaArgumentM8PsiBijective}
\FormulaThmDeltaKR[Argument (M8): Disjunktheit und Absorption der Kernfamilien]{%
 \begin{aligned}[t]
  &\text{(i)}\quad \mathcal C\subseteq\powerset(D)\\[-2pt]
  &\text{(ii)}\quad \mathcal C\cap\mathcal E=\varnothing\\[-2pt]
  &\text{(iii)}\quad \psi\colon\mathcal C\bij\mathcal C\cup\mathcal E
 \end{aligned}%
}{%
  \mathcal C\subseteq\powerset(D)\dsep
  \mathcal C\cap\mathcal E=\varnothing\dsep
  \psi\colon\mathcal C\bij\mathcal C\cup\mathcal E%
}{MogiljanskajaArgumentM8PsiBijective}{%
  \DeltaRow{Mengen}{D\dsep D'\dsep U\dsep V\dsep P\dsep c_2\dsep k
    \dsep\mathcal C\dsep\mathcal C_0\dsep\mathcal C_1
    \dsep\mathcal E\dsep H}
  \DeltaRow{Funktionen}{\sigma\dsep\theta\dsep
    \kappa_0\dsep\kappa_1\dsep\psi}
}
\MogProof{MogiljanskajaArgumentM8PsiBijective}{%
\begin{tabproofwide}
  \proofstepwidestar[]{V=U\cup\{c_2\}\dsep P\cup V=D'\subseteq D}{%
    \FormulaRefAuto{MogiljanskajaDPrimeDef}[def],
    \FormulaRefAuto{MogiljanskajaReserveMarkerFacts}[thm]}
  \proofstepwidestar[]{\mathcal C=\CoreFamily{P}{V}}{%
    \FormulaRefAuto{MogiljanskajaReserveMapsDef}[def]}
  \proofstepwidestar[1,2]{\mathcal C\subseteq\powerset(D')}{%
    \FormulaRefAuto{CoreFamilySubsetPowerset}[thm]}
  \proofstepwidestar[1,3]{\mathcal C\subseteq\powerset(D)}{%
    \text{\parbox[t]{\linewidth}{\raggedright \normalfont Monotonie der Potenzmenge}}}
  \proofstepwidestar[]{%
    \begin{aligned}[t]
      \mathcal C_0&=\CoreFamily{P}{U},&
      \mathcal C_1&=\CoreFamily{P\cup\{c_2\}}{U},\\[-2pt]
      \kappa_0&=\CoreTransport{P}{P}{\sigma},&
      \kappa_1&=\CoreTransport{P\cup\{c_2\}}{\{k\}}{\theta},\\[-2pt]
      \mathcal E&=\CoreFamily{\{k\}}{D},&
      \psi&=\CaseFun{\mathcal C_0}{\mathcal C_1}{\kappa_0}{\kappa_1}
    \end{aligned}}{%
    \FormulaRefAuto{MogiljanskajaReserveMapsDef}[def]}
  \proofstepwidestar[5]{H\in\mathcal E\vdash k\in H}{%
    \FormulaRefAuto{CoreFamilyMembership}[thm]}
  \proofstepwidestar[]{k\notin D}{%
    \FormulaRefAuto{MogiljanskajaArgumentM7AbsorptionHypotheses}[thm]}
  \proofstepwidestar[4,6,7]{\mathcal C\cap\mathcal E=\varnothing}{\text{\parbox[t]{\linewidth}{\raggedright %
    Kern \(k\) liegt außerhalb von \(D\)}}}
  \proofstepwidestar[]{%
    \begin{aligned}[t]
      &P\cap U=\varnothing\dsep c_2\notin P\cup U\dsep
        \{k\}\cap D=\varnothing,\\[-2pt]
      &\sigma\colon U\bij V\dsep\theta\colon U\bij D
    \end{aligned}}{%
    \FormulaRefAuto{MogiljanskajaArgumentM7AbsorptionHypotheses}[thm]}
  \proofstepwidestar[1,5,8,9]{\psi\colon\mathcal C\bij
    \mathcal C\cup\mathcal E}{%
    \FormulaRefAuto{CoreFamilyAbsorptionBijective}[thm]}
\end{tabproofwide}
}


\MogStatementTarget{MogiljanskajaArgumentM9EtaBijective}
\FormulaThmDeltaKR[Argument (M9): Zielzerlegung und Bijektivität von \(\eta\)]{%
 \begin{aligned}[t]
  &\text{(i)}\quad \powerset(D\cup\{k\}) =\mathcal O\cup\mathcal C\cup\mathcal E\\[-2pt]
  &\text{(ii)}\quad \mathcal O\cap\mathcal C=\mathcal O\cap\mathcal E =\mathcal C\cap\mathcal E=\varnothing\\[-2pt]
  &\text{(iii)}\quad \eta\colon\powerset(D)\bij\powerset(D\cup\{k\})
 \end{aligned}%
}{%
  \begin{aligned}[t]
    &\powerset(D\cup\{k\})
      =\mathcal O\cup\mathcal C\cup\mathcal E,\\[-2pt]
    &\mathcal O\cap\mathcal C=\mathcal O\cap\mathcal E
      =\mathcal C\cap\mathcal E=\varnothing,\\[-2pt]
    &\eta\colon\powerset(D)\bij\powerset(D\cup\{k\})
  \end{aligned}
}{MogiljanskajaArgumentM9EtaBijective}{%
  \DeltaRow{Mengen}{D\dsep P\dsep k\dsep\mathcal O\dsep\mathcal C
    \dsep\mathcal E\dsep H}
  \DeltaRow{Funktionen}{\psi\dsep\eta}
}
\MogProof{MogiljanskajaArgumentM9EtaBijective}{%
\begin{tabproofwide}
  \proofstepwidestarbreak[]{\mathcal E
    =\mathfrak I[\{k\},D\cup\{k\}]}{%
    \FormulaRefAuto{CoreFamilyBooleanInterval}[thm],
    \FormulaRefAuto{MogiljanskajaReserveMapsDef}[def]}
  \proofstepwidestar[]{k\notin D}{%
    \FormulaRefAuto{MogiljanskajaArgumentM7AbsorptionHypotheses}[thm]}
  \proofstepwidestar[]{\mathcal C\subseteq\powerset(D)}{%
    \FormulaRefAuto{MogiljanskajaArgumentM8PsiBijective}[thm]}
  \proofstepwidestar[]{P\neq\varnothing}{%
    \FormulaRefAuto{MogiljanskajaDPrimeDef}[def]}
  \proofstepwidestarbreak[4]{\varnothing\notin\mathcal C}{%
    \FormulaRefAuto{CoreFamilyMembership}[thm],
    \FormulaRefAuto{MogiljanskajaReserveMapsDef}[def]}
  \proofstepwidestar[]{\psi\colon\mathcal C\bij
    \mathcal C\cup\mathcal E}{%
    \FormulaRefAuto{MogiljanskajaArgumentM8PsiBijective}[thm]}
  \proofstepwidestarbreak[]{\mathcal O=\powerset(D)\setminus\mathcal C\dsep
    \eta=\CaseFun{\mathcal O}{\mathcal C}{\Id_{\mathcal O}}{\psi}}{%
    \FormulaRefAuto{MogiljanskajaReserveMapsDef}[def]}
  \proofstepwidestarbreak[1,2,3,5,6,7]{%
    \eta\colon\powerset(D)\bij\powerset(D\cup\{k\})}{%
    \FormulaRefAuto{ProtectedPowerSetExtensionBijective}[thm]}
  \proofstepwidestarbreak[1]{H\in\mathcal E\leftrightarrow
    H\subseteq D\cup\{k\}\land k\in H}{%
    \FormulaRefAuto{CoreFamilyMembership}[thm]}
  \proofstepwidestarbreak[]{H\subseteq D\cup\{k\}\dsep k\notin H
    \vdash H\subseteq D}{%
    \text{\parbox[t]{\linewidth}{\raggedright \normalfont Abspaltung des einzigen neuen Punktes}}}
  \proofstepwidestarbreak[9,10]{\powerset(D\cup\{k\})
    =\powerset(D)\cup\mathcal E}{\text{\parbox[t]{\linewidth}{\raggedright \normalfont Fall \ensuremath{k\in H} \normalfont{ oder } \ensuremath{k\notin H}}}}
  \proofstepwidestarbreak[]{\powerset(D)=\mathcal O\cup\mathcal C\dsep
    \mathcal O\cap\mathcal C=\varnothing}{%
    \FormulaRefAuto{MogiljanskajaReserveMapsDef}[def]}
  \proofstepwidestarbreak[11,12]{\powerset(D\cup\{k\})
    =\mathcal O\cup\mathcal C\cup\mathcal E}{\text{\parbox[t]{\linewidth}{\raggedright \normalfont Einsetzen}}}
  \proofstepwidestar[2,9,12]{\mathcal O\cap\mathcal E=\varnothing}{\text{\parbox[t]{\linewidth}{\raggedright %
    Mengen aus \(\mathcal O\) enthalten \(k\) nicht}}}
  \proofstepwidestar[]{\mathcal C\cap\mathcal E=\varnothing}{%
    \FormulaRefAuto{MogiljanskajaArgumentM8PsiBijective}[thm]}
  \proofstepwidestarbreak[12,14,15]{%
    \mathcal O\cap\mathcal C=\mathcal O\cap\mathcal E
      =\mathcal C\cap\mathcal E=\varnothing}{%
    \text{\parbox[t]{\linewidth}{\raggedright \normalfont paarweise Disjunktheit}}}
\end{tabproofwide}
}


\MogStatementTarget{MogiljanskajaArgumentM10EtaFixedPart}
\FormulaThmDeltaKR[Argument (M10): Nulltreue und Fixierung außerhalb der Reserve]{%
 \begin{aligned}[t]
  &\text{(i)}\quad \eta(\varnothing)=\varnothing\\[-2pt]
  &\text{(ii)}\quad \forall K\in\powerset(D)\, \bigl(K\notin\mathcal C\rightarrow\eta(K)=K\bigr)
 \end{aligned}%
}{%
  \eta(\varnothing)=\varnothing\dsep
  \forall K\in\powerset(D)\,
    \bigl(K\notin\mathcal C\rightarrow\eta(K)=K\bigr)%
}{MogiljanskajaArgumentM10EtaFixedPart}{%
  \DeltaRow{Mengen}{D\dsep P\dsep c_0\dsep\mathcal C
    \dsep\mathcal O\dsep K}
  \DeltaRow{Funktionen}{\eta}
}
\MogProof{MogiljanskajaArgumentM10EtaFixedPart}{%
\begin{tabproofwide}
  \proofstepwidestar[]{c_0\in P}{\FormulaRefAuto{MogiljanskajaDPrimeDef}[def]}
  \proofstepwidestar[]{P\neq\varnothing}{\FormulaRefAuto{a\in A\vdash A\neq\varnothing}[thm]{1}}
  \proofstepwidestar[]{\varnothing\notin\mathcal C}{%
    \FormulaRefAuto{CoreFamilyMembership}[thm]{2},
    \FormulaRefAuto{MogiljanskajaReserveMapsDef}[def]}
  \proofstepwidestar[]{\varnothing\in\mathcal O}{\FormulaRefAuto{MogiljanskajaReserveMapsDef}[def]{3}}
  \proofstepwidestar[]{\eta(\varnothing)=\varnothing}{\FormulaRefAuto{MogiljanskajaReserveMapsDef}[def]{4}}
  \proofstepwidestar[6]{K\in\powerset(D)}{\rA}
  \proofstepwidestar[7]{K\notin\mathcal C}{\rA}
  \proofstepwidestar[6,7]{K\in\mathcal O}{\FormulaRefAuto{MogiljanskajaReserveMapsDef}[def]{6,7}}
  \proofstepwidestar[6,7]{\eta(K)=K}{\FormulaRefAuto{MogiljanskajaReserveMapsDef}[def]{8}}
  \proofstepwidestarbreak[]{\forall K\in\powerset(D)\,
    (K\notin\mathcal C\rightarrow\eta(K)=K)}{\rUI{\rRI{6,\rRI{7,9}}}}
  \proofstepwidestarbreak[]{\eta(\varnothing)=\varnothing\land
    \forall K\in\powerset(D)\,
      (K\notin\mathcal C\rightarrow\eta(K)=K)}{\rAI{5,10}}
\end{tabproofwide}
}


\MogStatementTarget{MogiljanskajaArgumentM11RectangleProtected}
\FormulaThmDeltaK[Argument (M11): Kein Produktrechteck liegt in der Reserve]{%
  X,Y\in\powerset(T)\vdash R(X,Y)\notin\mathcal C%
}{MogiljanskajaArgumentM11RectangleProtected}{%
  \DeltaRow{Mengen}{L\dsep D'\dsep T\dsep P\dsep\mathcal C
    \dsep X\dsep Y\dsep R(X,Y)\dsep a_0\dsep a_1
    \dsep b_{00}\dsep b_{11}\dsep b_{01}}
}
\MogProof{MogiljanskajaArgumentM11RectangleProtected}{%
\begin{tabproofwide}
  \proofstepwidestar[1]{X,Y\in\powerset(T)}{\rA}
  \proofstepwidestar[2]{R(X,Y)\in\mathcal C}{\rA}
  \proofstepwidestarbreak[2]{P\subseteq R(X,Y)\dsep R(X,Y)\subseteq D'}{%
    \FormulaRefAuto{CoreFamilyMembership}[thm]{2},
    \FormulaRefAuto{MogiljanskajaReserveMapsDef}[def],
    \FormulaRefAuto{MogiljanskajaDPrimeDef}[def]}
  \proofstepwidestar[3]{b_{00},b_{11}\in R(X,Y)}{%
    \FormulaRefAuto{MogiljanskajaDPrimeDef}[def]{3}}
  \proofstepwidestar[4]{a_0\in X\cap L\dsep a_0\in Y\cap L}{%
    \FormulaRefAuto{MogiljanskajaProductPartsDef}[def],
    \FormulaRefAuto{MogiljanskajaBIndexUnique}[thm]}
  \proofstepwidestar[4]{a_1\in X\cap L\dsep a_1\in Y\cap L}{%
    \FormulaRefAuto{MogiljanskajaProductPartsDef}[def],
    \FormulaRefAuto{MogiljanskajaBIndexUnique}[thm]}
  \proofstepwidestar[5,6]{b_{01}\in R(X,Y)}{%
    \FormulaRefAuto{MogiljanskajaProductPartsDef}[def]}
  \proofstepwidestar[3,7]{b_{01}\in D'}{\rIE{7,3}}
  \proofstepwidestar[]{b_{01}\notin D'}{%
    \FormulaRefAuto{MogiljanskajaReserveMarkerFacts}[thm]}
  \proofstepwidestar[2]{\bot}{\rBI{8,9}}
  \proofstepwidestar[1]{R(X,Y)\notin\mathcal C}{\rCI{2,10}}
  \proofstepwidestar[]{X,Y\in\powerset(T)\vdash
    R(X,Y)\notin\mathcal C}{\rRI{1,11}}
\end{tabproofwide}
}


\Needspace{20\baselineskip}
\section{Globale Abbildung und Isomorphieschluss}

\MogStatementTarget{MogiljanskajaArgumentM12PhiBijective}
\FormulaThmDeltaKR[Argument (M12): Schichtzerlegung und globale Bijektion]{%
 \begin{aligned}[t]
  &\text{(i)}\quad A_*\cap D=\varnothing\\[-2pt]
  &\text{(ii)}\quad A_*\cap(D\cup\{k\})=\varnothing\\[-2pt]
  &\text{(iii)}\quad S=A_*\cup D\\[-2pt]
  &\text{(iv)}\quad T=A_*\cup(D\cup\{k\})\\[-2pt]
  &\text{(v)}\quad \Phi\colon\PowNE{S}\bij\PowNE{T}
 \end{aligned}%
}{%
  \begin{aligned}[t]
    &A_*\cap D=\varnothing\dsep
      A_*\cap(D\cup\{k\})=\varnothing,\\[-2pt]
    &S=A_*\cup D\dsep T=A_*\cup(D\cup\{k\}),\\[-2pt]
    &\Phi\colon\PowNE{S}\bij\PowNE{T}
  \end{aligned}
}{MogiljanskajaArgumentM12PhiBijective}{%
  \DeltaRow{Mengen}{A_*\dsep L\dsep D\dsep S\dsep T
    \dsep\mathbf0\dsep k}
  \DeltaRow{Funktionen}{\eta\dsep\Phi}
}
\MogProof{MogiljanskajaArgumentM12PhiBijective}{%
\begin{tabproofwide}
  \proofstepwidestar[]{A_*=L\cup\{\mathbf0\}}{%
    \FormulaRefAuto{MogiljanskajaPhiDef}[def]}
  \proofstepwidestar[1]{A_*\cap D=\varnothing}{%
    \FormulaRefAuto{MogiljanskajaArgumentM1WellDefined}[thm]}
  \proofstepwidestar[1]{A_*\cap(D\cup\{k\})=\varnothing}{%
    \FormulaRefAuto{MogiljanskajaArgumentM1WellDefined}[thm]}
  \proofstepwidestarbreak[1]{S=A_*\cup D\dsep
    T=A_*\cup(D\cup\{k\})}{%
    \FormulaRefAuto{MogiljanskajaPairDef}[def]}
  \proofstepwidestarbreak[]{\eta\colon\powerset(D)\bij
    \powerset(D\cup\{k\})}{%
    \FormulaRefAuto{MogiljanskajaArgumentM9EtaBijective}[thm]}
  \proofstepwidestar[]{\eta(\varnothing)=\varnothing}{%
    \FormulaRefAuto{MogiljanskajaArgumentM10EtaFixedPart}[thm]}
  \proofstepwidestar[2,3,4,5,6]{\Phi\colon\PowNE{S}\bij\PowNE{T}}{%
    \FormulaRefAuto{LayerPowerMapNonemptyBijective}[thm],
    \FormulaRefAuto{MogiljanskajaPhiDef}[def]}
\end{tabproofwide}
}


\MogStatementTarget{MogiljanskajaArgumentM13PhiInvariants}
\FormulaThmDeltaKR[Argument (M13): \(L\)-Invarianz von \(\Phi\)]{%
 \begin{aligned}[t]
  &\text{(i)}\quad X\in\PowNE{S}\vdash{} \Phi(X)=(X\cap A_*)\cup\eta(X\cap D)\\[-2pt]
  &\text{(ii)}\quad X\in\PowNE{S}\vdash{} \Phi(X)\cap L=X\cap L\\[-2pt]
  &\text{(iii)}\quad X\in\PowNE{S}\vdash{} \bigl(X\subseteq L\leftrightarrow\Phi(X)\subseteq L\bigr)
 \end{aligned}%
}{%
  \begin{aligned}[t]
    X\in\PowNE{S}\vdash{}&
      \Phi(X)=(X\cap A_*)\cup\eta(X\cap D),\\[-2pt]
    X\in\PowNE{S}\vdash{}&
      \Phi(X)\cap L=X\cap L,\\[-2pt]
    X\in\PowNE{S}\vdash{}&
      \bigl(X\subseteq L\leftrightarrow\Phi(X)\subseteq L\bigr)
  \end{aligned}
}{MogiljanskajaArgumentM13PhiInvariants}{%
  \DeltaRow{Mengen}{A_*\dsep L\dsep D\dsep S\dsep T\dsep X}
  \DeltaRow{Funktionen}{\eta\dsep\Phi}
}
\MogProof{MogiljanskajaArgumentM13PhiInvariants}{%
\begin{tabproofwide}
  \proofstepwidestarbreak[]{A_*\cap D=\varnothing\dsep
    A_*\cap(D\cup\{k\})=\varnothing\dsep
    S=A_*\cup D}{%
    \FormulaRefAuto{MogiljanskajaArgumentM12PhiBijective}[thm]}
  \proofstepwidestarbreak[]{\eta\colon\powerset(D)\bij
    \powerset(D\cup\{k\})}{%
    \FormulaRefAuto{MogiljanskajaArgumentM9EtaBijective}[thm]}
  \proofstepwidestar[]{\eta(\varnothing)=\varnothing}{%
    \FormulaRefAuto{MogiljanskajaArgumentM10EtaFixedPart}[thm]}
  \proofstepwidestar[1]{X\in\PowNE{S}\vdash X\subseteq A_*\cup D}{%
    \text{\parbox[t]{\linewidth}{\raggedright \normalfont Trägerzerlegung}}}
  \proofstepwidestarbreak[1,2,4]{\begin{aligned}[t]&X\in\PowNE{S}\vdash{}\\[-2pt]
    &\Phi(X)=(X\cap A_*)\cup\eta(X\cap D)\end{aligned}}{%
    \FormulaRefAuto{LayerPowerMapEvaluation}[thm],
    \FormulaRefAuto{MogiljanskajaPhiDef}[def]}
  \proofstepwidestarbreak[1,2,4]{X\in\PowNE{S}\vdash
    \Phi(X)\cap A_*=X\cap A_*}{%
    \FormulaRefAuto{LayerPowerMapStablePart}[thm]}
  \proofstepwidestar[]{L\subseteq A_*}{%
    \FormulaRefAuto{MogiljanskajaPhiDef}[def]}
  \proofstepwidestarbreak[6,7]{X\in\PowNE{S}\vdash
    \Phi(X)\cap L=X\cap L}{\text{\parbox[t]{\linewidth}{\raggedright \normalfont Schnitt mit der Teilschicht \ensuremath{L}}}}
  \proofstepwidestarbreak[1,2,3,4,7]{X\in\PowNE{S}\vdash
    \bigl(X\subseteq L\leftrightarrow\Phi(X)\subseteq L\bigr)}{%
    \FormulaRefAuto{LayerPowerMapPreservesSubsets}[thm]}
\end{tabproofwide}
}


\MogStatementTarget{MogiljanskajaArgumentM14ProductPartsInvariant}
\FormulaThmDeltaKR[Argument (M14): Invarianz der Produktanteile]{%
 \begin{aligned}[t]
  &\text{(i)}\quad X,Y\in\PowNE{S}\vdash{} R(\Phi(X),\Phi(Y))=R(X,Y)\\[-2pt]
  &\text{(ii)}\quad X,Y\in\PowNE{S}\vdash{} \varepsilon(\Phi(X),\Phi(Y))=\varepsilon(X,Y)\\[-2pt]
  &\text{(iii)}\quad X,Y\in\PowNE{S}\vdash{}\\[-2pt]
  &\qquad\Phi(X)\MogDiamond_{\mathcal P}\Phi(Y)
    =R(X,Y)\cup\varepsilon(X,Y)=X\MogStar_{\mathcal P}Y
 \end{aligned}%
}{%
  \begin{aligned}[t]
    X,Y\in\PowNE{S}\vdash{}&
      R(\Phi(X),\Phi(Y))=R(X,Y),\\[-2pt]
    X,Y\in\PowNE{S}\vdash{}&
      \varepsilon(\Phi(X),\Phi(Y))=\varepsilon(X,Y),\\[-2pt]
    X,Y\in\PowNE{S}\vdash{}&
      \Phi(X)\MogDiamond_{\mathcal P}\Phi(Y)
        =R(X,Y)\cup\varepsilon(X,Y)
        =X\MogStar_{\mathcal P}Y
  \end{aligned}
}{MogiljanskajaArgumentM14ProductPartsInvariant}{%
  \DeltaRow{Mengen}{L\dsep S\dsep T\dsep X\dsep Y
    \dsep R(X,Y)\dsep\varepsilon(X,Y)}
  \DeltaRow{Funktionen}{\Phi\dsep\MogStar_{\mathcal P}
    \dsep\MogDiamond_{\mathcal P}}
}
\MogProof{MogiljanskajaArgumentM14ProductPartsInvariant}{%
\begin{tabproofwide}
  \proofstepwidestarbreak[]{X,Y\in\PowNE{S}\vdash
    \Phi(X)\cap L=X\cap L\dsep
    \Phi(Y)\cap L=Y\cap L}{%
    \FormulaRefAuto{MogiljanskajaArgumentM13PhiInvariants}[thm]}
  \proofstepwidestarbreak[1]{X,Y\in\PowNE{S}\vdash
    R(\Phi(X),\Phi(Y))=R(X,Y)}{%
    \FormulaRefAuto{MogiljanskajaProductPartsDef}[def]}
  \proofstepwidestarbreak[]{%
    \begin{aligned}[t]
      X,Y\in\PowNE{S}&\vdash{}\\[-2pt]
      X\subseteq L&\leftrightarrow\Phi(X)\subseteq L,\\[-2pt]
      Y\subseteq L&\leftrightarrow\Phi(Y)\subseteq L
    \end{aligned}}{%
    \FormulaRefAuto{MogiljanskajaArgumentM13PhiInvariants}[thm]}
  \proofstepwidestarbreak[3]{X,Y\in\PowNE{S}\vdash
    \varepsilon(\Phi(X),\Phi(Y))=\varepsilon(X,Y)}{%
    \FormulaRefAuto{MogiljanskajaProductPartsDef}[def]}
  \proofstepwidestarbreak[]{X,Y\in\PowNE{S}\vdash
    \Phi(X),\Phi(Y)\in\PowNE{T}}{%
    \FormulaRefAuto{MogiljanskajaArgumentM12PhiBijective}[thm]}
  \proofstepwidestarbreak[2,4,5]{X,Y\in\PowNE{S}\vdash
    \Phi(X)\MogDiamond_{\mathcal P}\Phi(Y)
      =R(X,Y)\cup\varepsilon(X,Y)}{%
    \FormulaRefAuto{MogiljanskajaArgumentM6ProductClassification}[thm]}
  \proofstepwidestarbreak[]{X,Y\in\PowNE{S}\vdash
    X\MogStar_{\mathcal P}Y
      =R(X,Y)\cup\varepsilon(X,Y)}{%
    \FormulaRefAuto{MogiljanskajaArgumentM6ProductClassification}[thm]}
  \proofstepwidestarbreak[6,7]{X,Y\in\PowNE{S}\vdash
    \Phi(X)\MogDiamond_{\mathcal P}\Phi(Y)
      =R(X,Y)\cup\varepsilon(X,Y)
      =X\MogStar_{\mathcal P}Y}{%
    \text{\parbox[t]{\linewidth}{\raggedright \normalfont Gleichheitskette}}}
\end{tabproofwide}
}


\MogStatementTarget{MogiljanskajaArgumentM15ProductFixed}
\FormulaThmDeltaK[Argument (M15): Fixierung jedes ersten Mengenprodukts]{%
  X,Y\in\PowNE{S}\vdash
  \Phi(X\MogStar_{\mathcal P}Y)=X\MogStar_{\mathcal P}Y%
}{MogiljanskajaArgumentM15ProductFixed}{%
  \DeltaRow{Mengen}{A_*\dsep D\dsep S\dsep X\dsep Y\dsep Z
    \dsep R(X,Y)\dsep\varepsilon(X,Y)\dsep\mathcal C}
  \DeltaRow{Funktionen}{\eta\dsep\Phi\dsep\MogStar_{\mathcal P}}
}
\MogProof{MogiljanskajaArgumentM15ProductFixed}{%
\begin{tabproofwide}
  \proofstepwidestar[]{X,Y\in\PowNE{S}\dsep
    Z=X\MogStar_{\mathcal P}Y}{\rA}
  \proofstepwidestar[1]{X,Y\in\powerset(T)}{%
    \FormulaRefAuto{MogiljanskajaPairDef}[def]}
  \proofstepwidestar[1]{Z\in\PowNE{S}}{%
    \FormulaRefAuto{MogiljanskajaArgumentM2Semigroups}[thm],
    \FormulaRefAuto{NonemptyPowerSemigroup}[thm]}
  \proofstepwidestar[1]{Z=R(X,Y)\cup\varepsilon(X,Y)}{%
    \FormulaRefAuto{MogiljanskajaArgumentM6ProductClassification}[thm]{1}}
  \proofstepwidestar[]{R(X,Y)\subseteq D\dsep
    \varepsilon(X,Y)\subseteq\{\mathbf0\}\subseteq A_*}{%
    \FormulaRefAuto{MogiljanskajaProductPartsDef}[def],
    \FormulaRefAuto{MogiljanskajaPhiDef}[def]}
  \proofstepwidestar[4,5]{Z\cap D=R(X,Y)\dsep
    Z\cap A_*=\varepsilon(X,Y)}{%
    \FormulaRefAuto{MogiljanskajaArgumentM12PhiBijective}[thm]}
  \proofstepwidestar[2]{R(X,Y)\notin\mathcal C}{%
    \FormulaRefAuto{MogiljanskajaArgumentM11RectangleProtected}[thm]}
  \proofstepwidestar[5,7]{\eta(R(X,Y))=R(X,Y)}{%
    \FormulaRefAuto{MogiljanskajaArgumentM10EtaFixedPart}[thm]}
  \proofstepwidestar[3,6]{\Phi(Z)
    =(Z\cap A_*)\cup\eta(Z\cap D)}{%
    \FormulaRefAuto{MogiljanskajaArgumentM13PhiInvariants}[thm]}
  \proofstepwidestar[4,6,8,9]{\Phi(Z)
    =\varepsilon(X,Y)\cup R(X,Y)=Z}{\text{\parbox[t]{\linewidth}{\raggedright \normalfont Einsetzen}}}
  \proofstepwidestar[1,10]{\Phi(X\MogStar_{\mathcal P}Y)
    =X\MogStar_{\mathcal P}Y}{\text{\parbox[t]{\linewidth}{\raggedright %
    Definition von \(Z\)}}}
\end{tabproofwide}
}


\MogStatementTarget{MogiljanskajaArgumentM16PowerIsomorphism}
\FormulaThmDeltaK[Argument (M16): Direkter Homomorphie- und Isomorphieschluss]{%
  \begin{aligned}[t]
    &\SemigroupHom{\Phi}
      {\PowNE{S},\MogStar_{\mathcal P}}
      {\PowNE{T},\MogDiamond_{\mathcal P}}\land{}\\[-2pt]
    &\SemigroupIso{\Phi}
      {\PowNE{S},\MogStar_{\mathcal P}}
      {\PowNE{T},\MogDiamond_{\mathcal P}}
  \end{aligned}
}{MogiljanskajaArgumentM16PowerIsomorphism}{%
  \DeltaRow{Mengen}{S\dsep T\dsep X\dsep Y}
  \DeltaRow{Funktionen}{\Phi\dsep\MogStar_{\mathcal P}
    \dsep\MogDiamond_{\mathcal P}}
}
\MogProof{MogiljanskajaArgumentM16PowerIsomorphism}{%
\begin{tabproofwide}
  \proofstepwidestar[]{X,Y\in\PowNE{S}\vdash
    \Phi(X\MogStar_{\mathcal P}Y)
      =X\MogStar_{\mathcal P}Y}{%
    \FormulaRefAuto{MogiljanskajaArgumentM15ProductFixed}[thm]}
  \proofstepwidestar[]{X,Y\in\PowNE{S}\vdash
    \Phi(X)\MogDiamond_{\mathcal P}\Phi(Y)
      =X\MogStar_{\mathcal P}Y}{%
    \FormulaRefAuto{MogiljanskajaArgumentM14ProductPartsInvariant}[thm]}
  \proofstepwidestarbreak[1,2]{%
    \begin{aligned}[t]
      &X,Y\in\PowNE{S}\vdash{}\\[-2pt]
      &\quad\Phi(X\MogStar_{\mathcal P}Y)
        =\Phi(X)\MogDiamond_{\mathcal P}\Phi(Y)
    \end{aligned}}{%
    \text{\parbox[t]{\linewidth}{\raggedright \normalfont gemeinsame rechte Seite}}}
  \proofstepwidestar[]{%
    \SemigroupAlg{\PowNE{S}}{\MogStar_{\mathcal P}}\land
    \SemigroupAlg{\PowNE{T}}{\MogDiamond_{\mathcal P}}}{%
    \FormulaRefAuto{MogiljanskajaArgumentM2Semigroups}[thm],
    \FormulaRefAuto{NonemptyPowerSemigroup}[thm]}
  \proofstepwidestar[]{\Phi\colon\PowNE{S}\bij\PowNE{T}}{%
    \FormulaRefAuto{MogiljanskajaArgumentM12PhiBijective}[thm]}
  \proofstepwidestar[3,4,5]{\SemigroupHom{\Phi}
    {\PowNE{S},\MogStar_{\mathcal P}}
    {\PowNE{T},\MogDiamond_{\mathcal P}}}{%
    \FormulaRefAuto{SemigroupHomDef}[def]}
  \proofstepwidestar[5,6]{\SemigroupIso{\Phi}
    {\PowNE{S},\MogStar_{\mathcal P}}
    {\PowNE{T},\MogDiamond_{\mathcal P}}}{%
    \FormulaRefAuto{SemigroupIsoDef}[def]}
\end{tabproofwide}
}


\MogStatementTarget{MogiljanskajaArgumentM17CounterexampleWitness}
\FormulaThmDeltaK[Argument (M17): Expliziter Schlusszeuge]{%
  \begin{aligned}[t]
    &\SemigroupAlg{S}{\MogStar}\land
      \SemigroupAlg{T}{\MogDiamond}\land{}\\[-2pt]
    &\neg\Finite{S}\land\neg\Finite{T}\land{}\\[-2pt]
    &\neg\exists f\,
      \SemigroupIso{f}{S,\MogStar}{T,\MogDiamond}\land{}\\[-2pt]
    &\SemigroupIso{\Phi}
      {\PowNE{S},\MogStar_{\mathcal P}}
      {\PowNE{T},\MogDiamond_{\mathcal P}}
  \end{aligned}
}{MogiljanskajaArgumentM17CounterexampleWitness}{%
  \DeltaRow{Mengen}{S\dsep T}
  \DeltaRow{Funktionen}{f\dsep\Phi\dsep\MogStar\dsep\MogDiamond
    \dsep\MogStar_{\mathcal P}\dsep\MogDiamond_{\mathcal P}}
}
\MogProof{MogiljanskajaArgumentM17CounterexampleWitness}{%
\begin{tabproofwide}
  \proofstepwidestar[]{\SemigroupAlg{S}{\MogStar}\land
    \SemigroupAlg{T}{\MogDiamond}}{%
    \FormulaRefAuto{MogiljanskajaArgumentM2Semigroups}[thm]}
  \proofstepwidestar[]{\neg\Finite{S}\land\neg\Finite{T}}{%
    \FormulaRefAuto{MogiljanskajaArgumentM3Infinite}[thm]}
  \proofstepwidestar[]{\neg\exists f\,
    \SemigroupIso{f}{S,\MogStar}{T,\MogDiamond}}{%
    \FormulaRefAuto{MogiljanskajaArgumentM5NotIsomorphic}[thm]}
  \proofstepwidestar[]{\SemigroupIso{\Phi}
    {\PowNE{S},\MogStar_{\mathcal P}}
    {\PowNE{T},\MogDiamond_{\mathcal P}}}{%
    \FormulaRefAuto{MogiljanskajaArgumentM16PowerIsomorphism}[thm]}
  \proofstepwidestar[1,2,3,4]{%
    \begin{aligned}[t]
      &\SemigroupAlg{S}{\MogStar}\land
        \SemigroupAlg{T}{\MogDiamond}\land{}\\[-2pt]
      &\neg\Finite{S}\land\neg\Finite{T}\land{}\\[-2pt]
      &\neg\exists f\,
        \SemigroupIso{f}{S,\MogStar}{T,\MogDiamond}\land{}\\[-2pt]
      &\SemigroupIso{\Phi}
        {\PowNE{S},\MogStar_{\mathcal P}}
        {\PowNE{T},\MogDiamond_{\mathcal P}}
    \end{aligned}}{\text{\parbox[t]{\linewidth}{\raggedright \normalfont Bündelung}}}
\end{tabproofwide}
}


